% ===== PRÉAMBULE CANONIQUE — à coller VERBATIM en tête de CHAQUE .tex =====
\documentclass[11pt,a4paper]{article}
\usepackage[utf8]{inputenc}\usepackage[T1]{fontenc}\usepackage{lmodern}
\usepackage[french]{babel}
\usepackage{amsmath,amssymb,mathtools}\usepackage{icomma}
\usepackage[version=4]{mhchem}          % \ce{...} : équations & formules
\usepackage{siunitx}\sisetup{locale=FR} % \qty{}{}, \si{}, \ang{}
\usepackage{chemfig}                    % \chemfig{...} : structures développées
\usepackage{xcolor}\usepackage{colortbl}
\usepackage{tikz}\usepackage{pgfplots}\pgfplotsset{compat=1.18}
\usepackage[most]{tcolorbox}\usepackage{enumitem}\usepackage{array,booktabs}
\usepackage[a4paper,margin=2.2cm]{geometry}
\usetikzlibrary{arrows.meta,calc,angles,quotes,positioning,patterns,backgrounds,shapes.geometric}
\definecolor{primary}{RGB}{222,49,99}\definecolor{primarydark}{RGB}{150,22,55}
\definecolor{defcol}{RGB}{0,114,178}\definecolor{methcol}{RGB}{52,92,148}
\definecolor{excol}{RGB}{0,135,90}\definecolor{attncol}{RGB}{200,40,55}
\tcbset{boxbase/.style={enhanced,breakable,boxrule=0.9pt,arc=2.5pt,
  left=3mm,right=3mm,top=2.3mm,bottom=2.3mm,fonttitle=\bfseries,coltitle=white}}
% --- boîtes TUTORIEL (noms EXACTS, ne pas renommer) ---
\newtcolorbox{cle}[1][]{boxbase,colback=defcol!6,colframe=defcol,
  title={\textsc{À retenir}\ifx\relax#1\relax\relax\else\textmd{ : \itshape #1}\fi}}
\newtcolorbox{methode}[1][]{boxbase,colback=methcol!7,colframe=methcol,sharp corners,
  title={\textsc{Méthode}\ifx\relax#1\relax\relax\else\textmd{ --- \itshape #1}\fi}}
\newtcolorbox{exemplebox}[1][]{boxbase,colback=excol!5,colframe=excol,
  title={\textsc{Exemple}\ifx\relax#1\relax\relax\else\textmd{ : \itshape #1}\fi}}
\newtcolorbox{piege}[1][]{boxbase,colback=attncol!6,colframe=attncol,
  title={\textsc{Piège}\ifx\relax#1\relax\relax\else\textmd{ : \itshape #1}\fi}}
\newtcolorbox{remarque}[1][]{boxbase,colback=black!4,colframe=black!45,
  title={\textsc{Remarque}\ifx\relax#1\relax\relax\else\textmd{ : \itshape #1}\fi}}
% --- boîtes EXERCICE / CORRIGÉ (noms EXACTS, auto-comptées) ---
\newtcolorbox[auto counter]{exo}[1][]{boxbase,colback=primary!4,colframe=primary,
  title={\textsc{Exercice}~\thetcbcounter\ifx\relax#1\relax\relax\else\textmd{ --- \itshape #1}\fi}}
\newtcolorbox[auto counter]{corrige}[1][]{boxbase,colback=excol!4,colframe=excol,
  title={\textsc{Corrigé}~\thetcbcounter\ifx\relax#1\relax\relax\else\textmd{ --- \itshape #1}\fi}}
\newcommand{\R}{\mathbb{R}}\newcommand{\N}{\mathbb{N}}\newcommand{\Z}{\mathbb{Z}}\newcommand{\ds}{\displaystyle}
% (un \usepackage{fancyhdr}+en-tête est possible ; il est ignoré côté web)
% =========================================================================
\begin{document}
\sisetup{per-mode=symbol}
\begin{exo}[addition et soustraction : les décimales]
Effectue, puis arrondis au bon nombre de \emph{décimales} (celui de la donnée la moins précise).
\begin{enumerate}[label=\alph*)]
  \item $123,25 + 46,0 + 86,257$
  \item $8,11 - 8,0$
  \item $12,4 + 1,07$
  \item $45,6 - 2,17$
\end{enumerate}
\end{exo}

\begin{exo}[multiplication et division : les CS]
Effectue, puis arrondis au bon nombre de \emph{chiffres significatifs} (celui de la donnée la moins riche).
\begin{enumerate}[label=\alph*)]
  \item $2,5 \times 3,42$
  \item $0,0361 \div 12,0$
  \item $4,00 \times 2,1$
  \item $6,022 \div 2,0$
\end{enumerate}
\end{exo}

\begin{exo}[quelle règle, quel rang ?]
\emph{Sans} faire le calcul exact, indique la règle à appliquer (décimales ou CS) et le nombre de rangs
à conserver dans le résultat.
\begin{enumerate}[label=\alph*)]
  \item $3,14 \times 2,5$
  \item $100,2 + 5,67$
  \item $8,0 \div 1,333$
  \item $12,00 - 3,1$
\end{enumerate}
\end{exo}

\begin{exo}[aire et volume]
Donne le résultat avec le bon nombre de chiffres significatifs.
\begin{enumerate}[label=\alph*)]
  \item l'aire d'un rectangle de côtés $\qty{2,5}{\centi\metre}$ et $\qty{4,0}{\centi\metre}$
  \item le volume d'un cube d'arête $\qty{3,0}{\centi\metre}$
\end{enumerate}
\end{exo}

\end{document}
